\ifdefined\gkver\else\def\gkver{student}\fi
\documentclass[\gkver]{gaokaozhenti}
\begin{document}

\chapter{2008年宁夏理综}
%% paperType: 理综物理部分

\begin{enumerate}
\item
%% number: 14
%% paperName: 2008年宁夏理综第 14 题 6 分
%% typeId: 1
%% type: 单选题
%% chapter: 磁场
%% point: 左手定则
%% method: 无
%% score: 6
%% degree: 650
%% duplicateId: 0
%% body:
在等边三角形的三个顶点 a、b、c 处，各有一条长直导线垂直穿过纸面，导线中通有大小相等的恒定电流，方向如图。过 c 点的导线所受安培力的方向\xzanswer{C}

\onepicture[w=4.5cm]{figs/14.png}

\fourchoices[answer=C]
{与 ab 边平行，竖直向上}
{与 ab 边平行，竖直向下}
{与 ab 边垂直，指向左边}
{与 ab 边垂直，指向右边}

%% answer: C
%% memo:
\memoanswer{
导线 a 在 c 处产生的磁场方向由安培定则可判断，即垂直 ac 向左；同理导线 b 在 c 处产生的磁场方向垂直 bc 向下。由平行四边形定则，过 c 点的合磁场方向平行于 ab，根据左手定则可判断导线 c 受到的安培力垂直 ab 边，指向左边。
}

\item
%% number: 15
%% paperName: 2008年宁夏理综第 15 题 6 分
%% typeId: 2
%% type: 多选题
%% chapter: 电路
%% point: 欧姆定律
%% method: 无
%% score: 6
%% degree: 650
%% duplicateId: 0
%% body:
一个 T 型电路如图所示，电路中的电阻 $R_{1}=10\UO$，$R_{2}=120\UO$，$R_{3}=40\UO$。另有一测试电源，电动势为 $100\UV$，内阻忽略不计。则\xzanswer{AC}

\onepicture[w=5.0cm]{figs/15.png}

\fourchoices[answer=AC]
{当 cd 端短路时，ab 之间的等效电阻是 $40\UO$}
{当 ab 端短路时，cd 之间的等效电阻是 $40\UO$}
{当 ab 两端接通测试电源时，cd 两端的电压为 $80\UV$}
{当 cd 两端接通测试电源时，ab 两端的电压为 $80\UV$}

%% answer: AC
%% memo:
\memoanswer{
当 cd 端短路时，$R_{2}$ 与 $R_{3}$ 并联电阻为 $30\UO$，后与 $R_{1}$ 串联，ab 间等效电阻为 $40\UO$，A 对；若 ab 端短路时，$R_{1}$ 与 $R_{3}$ 并联电阻为 $8\UO$，后与 $R_{2}$ 串联，cd 间等效电阻为 $128\UO$，B 错；ab 两端接通测试电源时，电阻 $R_{2}$ 未接入电路，cd 两端的电压即为 $R_{3}$ 的电压，$U_{cd}=\frac{40}{50}\times100\UV=80\UV$，C 对；cd 两端接通测试电源时，电阻 $R_{1}$ 未接入电路，ab 两端电压即为 $R_{3}$ 的电压，$U_{ab}=\frac{40}{160}\times100\UV=25\UV$，D 错。
}

\item
%% number: 16
%% paperName: 2008年宁夏理综第 16 题 6 分
%% typeId: 1
%% type: 单选题
%% chapter: 电磁感应
%% point: 右手定则
%% method: 无
%% score: 6
%% degree: 650
%% duplicateId: 0
%% body:
如图所示，同一平面内的三条平行导线串有两个电阻 $R$ 和 $r$，导体棒 PQ 与三条导线接触良好；匀强磁场的方向垂直纸面向里。导体棒的电阻可忽略。当导体棒向左滑动时，下列说法正确的是\xzanswer{B}

\onepicture[w=5.0cm]{figs/16.png}

\fourchoices[answer=B]
{流过 $R$ 的电流为由 d 到 c，流过 $r$ 的电流为由 b 到 a}
{流过 $R$ 的电流为由 c 到 d，流过 $r$ 的电流为由 b 到 a}
{流过 $R$ 的电流为由 d 到 c，流过 $r$ 的电流为由 a 到 b}
{流过 $R$ 的电流为由 c 到 d，流过 $r$ 的电流为由 a 到 b}

%% answer: B
%% memo:
\memoanswer{
根据右手定则可判断 PQ 作为电源，Q 端电势高。在 PQcd 回路中电流为逆时针方向，即流过 $R$ 的电流为由 c 到 d；在电阻 $r$ 的回路中电流为顺时针方向，即流过 $r$ 的电流为由 b 到 a。
}

\item
%% number: 17
%% paperName: 2008年宁夏理综第 17 题 6 分
%% typeId: 1
%% type: 单选题
%% chapter: 直线运动
%% point: 相遇问题
%% method: 无
%% score: 6
%% degree: 650
%% duplicateId: 0
%% body:
甲乙两车在公路上沿同一方向做直线运动，它们的 $v-t$ 图象如图所示。两图象在 $t=t_{1}$ 时相交于 P 点，P 在横轴上的投影为 Q，$\Delta OPQ$ 的面积为 $S$。在 $t=0$ 时刻，乙车在甲车前面，相距为 $d$。已知此后两车相遇两次，且第一次相遇的时刻为 $t'$，则下面四组 $t'$ 和 $d$ 的组合可能是\xzanswer{D}

\onepicture[w=5.0cm]{figs/17.png}

\fourchoices[answer=D]
{$t'=t_{1}$，$d=S$}
{$t'=\frac{1}{2}t_{1}$，$d=\frac{1}{4}S$}
{$t'=\frac{1}{2}t_{1}$，$d=\frac{1}{2}S$}
{$t'=\frac{1}{2}t_{1}$，$d=\frac{3}{4}S$}

%% answer: D
%% memo:
\memoanswer{
在 $t_{1}$ 时刻如果甲车没有追上乙车，以后就不可能追上了，故 $t'<t_{1}$，A 错；从图象看，甲在 $t_{1}$ 时间内的位移比乙的多 $S$，当 $t'=\frac{1}{2}t_{1}$ 时，甲的面积比乙的面积多出 $\frac{3}{4}S$，即相距 $d=\frac{3}{4}S$，D 正确。此类问题要抓住图象交点的物理意义，过了这个时刻不能相遇，以后就不可能相遇。
}

\item
%% number: 18
%% paperName: 2008年宁夏理综第 18 题 6 分
%% typeId: 1
%% type: 单选题
%% chapter: 功和能
%% point: 功
%% method: 无
%% score: 6
%% degree: 650
%% duplicateId: 0
%% body:
一滑块在水平地面上沿直线滑行，$t=0$ 时其速度为 $1\Ums$。从此刻开始滑块运动方向上再施加一水平作用力 $F$，力 $F$ 和滑块的速度 $v$ 随时间的变化规律分别如图 a 和图 b 所示。设在第 1 秒内、第 2 秒内、第 3 秒内力 $F$ 对滑块做的功分别为 $W_{1}$、$W_{2}$、$W_{3}$，则以下关系正确的是\xzanswer{B}

\onepicture[w=7.5cm]{figs/18.png}

\fourchoices[answer=B]
{$W_{1}=W_{2}=W_{3}$}
{$W_{1}<W_{2}<W_{3}$}
{$W_{1}<W_{3}<W_{2}$}
{$W_{1}=W_{2}<W_{3}$}

%% answer: B
%% memo:
\memoanswer{
力 $F$ 做的功等于每段恒力 $F$ 与该段滑块运动位移之积（位移为 $v-t$ 图象中图线与坐标轴围成的面积）。第 1 秒内 $F=1\UN$，位移为一个小三角形面积 $S=0.5\Um$，$W_{1}=0.5\UJ$；第 2 秒内 $F=3\UN$，位移也为 $S$，$W_{2}=1.5\UJ$；第 3 秒内 $F=2\UN$，位移为两个小三角形面积 $2S=1\Um$，$W_{3}=2\UJ$。故 $W_{1}<W_{2}<W_{3}$，B 正确。
}

\item
%% number: 19
%% paperName: 2008年宁夏理综第 19 题 6 分
%% typeId: 1
%% type: 单选题
%% chapter: 交流电
%% point: 交流电
%% method: 无
%% score: 6
%% degree: 650
%% duplicateId: 0
%% body:
如图 a 所示，一矩形线圈 abcd 放置在匀强磁场中，并绕过 ab、cd 中点的轴 OO$'$ 以角速度 $\omega$ 逆时针匀速转动。若以线圈平面与磁场夹角 $\theta=45^{\circ}$ 时（如图 b）为计时起点，并规定当电流自 a 流向 b 时电流方向为正。则下列四幅图中正确的是\xzanswer{D}

\onepicture[w=6.0cm]{figs/19.png}

\fourchoices[answer=D, ispicture=true, h=1.8cm]
{figs/19a.png}
{figs/19b.png}
{figs/19c.png}
{figs/19d.png}

%% answer: D
%% memo:
\memoanswer{
从图 a 可看出线圈从垂直于中性面开始旋转，由楞次定律可判断初始时刻电流方向为 b 到 a，故瞬时电流的表达式为 $i=-i_{\text{m}}\cos\left(\frac{\pi}{4}+\omega t\right)$，则图象为 D 所描述。
}

\item
%% number: 20
%% paperName: 2008年宁夏理综第 20 题 6 分
%% typeId: 2
%% type: 多选题
%% chapter: 运动和力
%% point: 牛顿第二定律
%% method: 临界
%% score: 6
%% degree: 650
%% duplicateId: 0
%% body:
一有固定斜面的小车在水平面上做直线运动，小球通过细绳与车顶相连。小球某时刻正处于图示状态。设斜面对小球的支持力为 $N$，细绳对小球的拉力为 $T$，关于此时刻小球的受力情况，下列说法正确的是\xzanswer{AB}

\onepicture[w=4.5cm]{figs/20.png}

\fourchoices[answer=AB]
{若小车向左运动，$N$ 可能为零}
{若小车向左运动，$T$ 可能为零}
{若小车向右运动，$N$ 不可能为零}
{若小车向右运动，$T$ 不可能为零}

%% answer: AB
%% memo:
\memoanswer{
对小球受力分析：当 $N$ 为零时，小球合外力水平向右，加速度向右，故小车可能向右加速运动或向左减速运动，A 对 C 错；当 $T$ 为零时，小球合外力水平向左，加速度向左，故小车可能向右减速运动或向左加速运动，B 对 D 错。解题时抓住 $N$、$T$ 为零时受力分析的临界条件，小球与车相对静止，说明小球和小车只能有水平的加速度。
}

\item
%% number: 21
%% paperName: 2008年宁夏理综第 21 题 6 分
%% typeId: 1
%% type: 单选题
%% chapter: 电路
%% point: 电容
%% method: 无
%% score: 6
%% degree: 450
%% duplicateId: 0
%% body:
如图所示，$C$ 为中间插有电介质的电容器，a 和 b 为其两极板；a 板接地；P 和 Q 为两竖直放置的平行金属板，在两板间用绝缘线悬挂一带电小球；P 板与 b 板用导线相连，Q 板接地。开始时悬线静止在竖直方向，在 b 板带电后，悬线偏转了角度 $\alpha$。在以下方法中，能使悬线的偏角 $\alpha$ 变大的是\xzanswer{BC}

\onepicture[w=5.0cm]{figs/21.png}

\fourchoices[answer=BC]
{缩小 a、b 间的距离}
{加大 a、b 间的距离}
{取出 a、b 两极板间的电介质}
{换一块形状大小相同、介电常数更大的电介质}

%% answer: BC
%% memo:
\memoanswer{
a 板与 Q 板电势恒为零，b 板和 P 板电势总相同，故两个电容器的电压相等，且两板电荷量 $q$ 视为不变。要使悬线的偏角增大，即电压 $U$ 增大，也就是要减小电容器的电容 $C$。对电容器 $C$，由 $C=\frac{q}{U}=\frac{\varepsilon S}{4\pi kd}$ 可知，可以通过增大板间距 $d$、减小介电常数 $\varepsilon$、减小极板正对面积 $S$ 来实现，故 B、C 正确。
}

\item
%% number: 22
%% paperName: 2008年宁夏理综第 22 题 9 分
%% typeId: 4
%% type: 实验题
%% chapter: 直线运动
%% point: 打点计时器公式
%% method: 无
%% score: 9
%% degree: 650
%% duplicateId: 0
%% body:
\begin{enumerate}
	\item
	Ⅰ．右图为一正在测量中的多用电表表盘。
	
	\onepicture[w=5.5cm, align=right]{figs/22.png}
	
	（1）如果是用 $\times10\UO$ 档测量电阻，则读数为 \tkanswer[1.5cm]{} $\UO$。
	
	（2）如果是用直流 $10\UmA$ 档测量电流，则读数为 \tkanswer[1.5cm]{} $\UmA$。
	
	（3）如果是用直流 $5\UV$ 档测量电压，则读数为 \tkanswer[1.5cm]{} $\UV$。
	
	\item
	Ⅱ．物理小组在一次探究活动中测量滑块与木板之间的动摩擦因数。实验装置如图，一表面粗糙的木板固定在水平桌面上，一端装有定滑轮；木板上有一滑块，其一端与电磁打点计时器的纸带相连，另一端通过跨过定滑轮的细线与托盘连接。打点计时器使用的交流电源的频率为 $50\UHz$。开始实验时，在托盘中放入适量砝码，滑块开始做匀加速运动，在纸带上打出一系列小点。
	
	\onepicture[w=5.0cm, align=right]{figs/22a.png}
	
	（1）上图给出的是实验中获取的一条纸带的一部分：0、1、2、3、4、5、6、7 是计数点，每相邻两计数点间还有 4 个打点（图中未标出），计数点间的距离如图所示。根据图中数据计算的加速度 $a=$ \tkanswer[2.0cm]{} $\Umsq$（保留三位有效数字）。
	
	\onepicture[w=9.0cm, align=right]{figs/22b.png}
	
	（2）回答下列两个问题：
	
	①为测量动摩擦因数，下列物理量中还应测量的有 \tkanswer[1.0cm]{}。（填入所选物理量前的字母）
	
	（A）木板的长度 $l$　（B）木板的质量 $m_{1}$　（C）滑块的质量 $m_{2}$　（D）托盘和砝码的总质量 $m_{3}$　（E）滑块运动的时间 $t$
	
	②测量①中所选定的物理量时需要的实验器材是 \tkanswer[2.0cm]{}。
	
	（3）滑块与木板间的动摩擦因数 $\mu=$ \tkanswer[3.0cm]{}（用被测物理量的字母表示，重力加速度为 $g$）。与真实值相比，测量的动摩擦因数 \tkanswer[1.5cm]{}（填“偏大”或“偏小”）。写出支持你的看法的一个论据：\tkanswer[4.0cm]{}。
\end{enumerate}

%% answer:
\jdanswer{
\begin{enumerate}
	\item
	Ⅰ．（1）$60\UO$；（2）$7.18\UmA$；（3）$3.59\UV$。
	\item
	Ⅱ．（1）$0.495\sim0.497\Umsq$（或 $0.496\Umsq$）；（2）①CD；②天平；（3）$\mu=\frac{m_{3}g-(m_{2}+m_{3})a}{m_{2}g}$，偏大，因滑轮与纸带存在阻力，使测得的动摩擦因数偏大。
\end{enumerate}
}

%% memo:
\memoanswer{
Ⅰ．欧姆档在最上面的一排数据读取，读数为 $6\times10\UO=60\UO$；电流档测量读取中间三排数据的最底下一排数据，读数为 $7.18\UmA$；同样直流电压档测量读取中间三排数据的中间一排数据较好，读数为 $35.9\times0.1\UV=3.59\UV$。

Ⅱ．对纸带的研究直接利用逐差法取平均值计算加速度。对滑块与托盘组成的系统，由牛顿第二定律有 $m_{3}g-\mu m_{2}g=(m_{2}+m_{3})a$，故 $\mu=\frac{m_{3}g-(m_{2}+m_{3})a}{m_{2}g}$；由于滑轮与纸带间存在阻力，测得的动摩擦因数偏大。
}

\item
%% number: 23
%% paperName: 2008年宁夏理综第 23 题 15 分
%% typeId: 6
%% type: 计算题
%% chapter: 曲线运动
%% point: 双星问题
%% method: 无
%% score: 15
%% degree: 450
%% duplicateId: 0
%% body:
天文学家将相距较近、仅在彼此的引力作用下运行的两颗恒星称为双星。双星系统在银河系中很普遍。利用双星系统中两颗恒星的运动特征可推算出它们的总质量。已知某双星系统中两颗恒星围绕它们连线上的某一固定点分别做匀速圆周运动，周期均为 $T$，两颗恒星之间的距离为 $r$，试推算这个双星系统的总质量。（引力常量为 $G$）

%% answer:
\jdanswer{
$m_{1}+m_{2}=\frac{4\pi^{2}r^{3}}{GT^{2}}$
}

%% memo:
\memoanswer{
设两颗恒星的质量分别为 $m_{1}$、$m_{2}$，做圆周运动的半径分别为 $r_{1}$、$r_{2}$，角速度分别为 $\omega_{1}$、$\omega_{2}$。根据题意有 $\omega_{1}=\omega_{2}$（①），$r_{1}+r_{2}=r$（②）。

根据万有引力定律和牛顿定律，有 $G\frac{m_{1}m_{2}}{r^{2}}=m_{1}\omega_{1}^{2}r_{1}$（③），$G\frac{m_{1}m_{2}}{r^{2}}=m_{2}\omega_{2}^{2}r_{2}$（④）。

联立以上各式解得 $r_{1}=\frac{m_{2}r}{m_{1}+m_{2}}$（⑤）。根据角速度与周期的关系知 $\omega_{1}=\omega_{2}=\frac{2\pi}{T}$（⑥）。

联立③⑤⑥式解得 $m_{1}+m_{2}=\frac{4\pi^{2}r^{3}}{GT^{2}}$（⑦）。
}

\item
%% number: 24
%% paperName: 2008年宁夏理综第 24 题 17 分
%% typeId: 6
%% type: 计算题
%% chapter: 磁场
%% point: 带电粒子在磁场中的运动
%% method: 无
%% score: 17
%% degree: 450
%% duplicateId: 0
%% body:
如图所示，在 $xOy$ 平面的第一象限有一匀强电场，电场的方向平行于 $y$ 轴向下；在 $x$ 轴和第四象限的射线 OC 之间有一匀强磁场，磁感应强度的大小为 $B$，方向垂直于纸面向外。有一质量为 $m$、带有电荷量 $+q$ 的质点由电场左侧平行于 $x$ 轴射入电场。质点到达 $x$ 轴上 A 点时，速度方向与 $x$ 轴的夹角为 $\varphi$，A 点与原点 O 的距离为 $d$。接着，质点进入磁场，并垂直于 OC 飞离磁场。不计重力影响。若 OC 与 $x$ 轴的夹角为 $\varphi$，求

\begin{enumerate}
	\item
	粒子在磁场中运动速度的大小；
	\item
	匀强电场的场强大小。
\end{enumerate}

\onepicture[w=5.0cm, align=right]{figs/24.png}

%% answer:
\jdanswer{
\begin{enumerate}
	\item
	$v=\frac{qBd\sin\varphi}{m}$
	\item
	$E=\frac{qB^{2}d\sin^{3}\varphi\cos\varphi}{m}$
\end{enumerate}
}

%% memo:
\memoanswer{
（1）质点在磁场中的轨迹为一圆弧。由于质点飞离磁场时速度垂直于 OC，故圆弧的圆心在 OC 上。依题意，质点轨迹与 $x$ 轴的交点为 A，过 A 点作与 A 点速度方向垂直的直线，与 OC 交于 $O'$。由几何关系知 $AO'\perp OC$，$O'$ 是圆弧的圆心。设圆弧的半径为 $R$，则有 $R=d\sin\varphi$（①）。

\onepicture[w=4.5cm]{figs/24_an.png}

由洛伦兹力公式和牛顿第二定律得 $qvB=m\frac{v^{2}}{R}$（②）。将①式代入②式得 $v=\frac{qBd\sin\varphi}{m}$（③）。

（2）质点在电场中的运动为类平抛运动。设质点射入电场的速度为 $v_{0}$，在电场中的加速度为 $a$，运动时间为 $t$，则有 $v_{0}=v\cos\varphi$（④），$v\sin\varphi=at$（⑤），$d=v_{0}t$（⑥）。

联立④⑤⑥得 $a=\frac{v^{2}\sin\varphi\cos\varphi}{d}$（⑦）。设电场强度的大小为 $E$，由牛顿第二定律得 $qE=ma$（⑧）。联立③⑦⑧得 $E=\frac{qB^{2}d\sin^{3}\varphi\cos\varphi}{m}$（⑨）。
}

\item
%% number: 30
%% paperName: 2008年宁夏理综第 30 题 10 分
%% typeId: 6
%% type: 计算题
%% chapter: 力物体的平衡
%% point: 力矩平衡
%% method: 无
%% score: 10
%% degree: 650
%% duplicateId: 0
%% body:
\begin{enumerate}
	\item
	（5 分）图示为某一皮带传动装置。主动轮的半径为 $r_{1}$，从动轮的半径为 $r_{2}$。已知主动轮做顺时针转动，转速为 $n$，转动过程中皮带不打滑。下列说法正确的是 \tkanswer[1.0cm]{}。（填入选项前的字母，有填错的不得分）
	
	\onepicture[w=4.5cm]{figs/30a.png}
	
	（A）从动轮做顺时针转动
	
	（B）从动轮做逆时针转动
	
	（C）从动轮的转速为 $\frac{nr_{1}}{r_{2}}$
	
	（D）从动轮的转速为 $\frac{nr_{2}}{r_{1}}$
	
	\item
	（10 分）一足够长的斜面，最高点为 O 点。有一长为 $l=1.00\Um$ 的木条 AB，A 端在斜面上，B 端伸出斜面外。斜面与木条间的摩擦力足够大，以致木条不会在斜面上滑动。在木条 A 端固定一个质量为 $M=2.00\Ukg$ 的重物（可视为质点），B 端悬挂一个质量为 $m=0.50\Ukg$ 的重物。若要使木条不脱离斜面，在下列两种情况下，OA 的长度各需满足什么条件？
	
	\onepicture[w=5.5cm, align=right]{figs/30b.png}
	
	（Ⅰ）木条的质量可以忽略不计。
	
	（Ⅱ）木条质量为 $m'=0.50\Ukg$，分布均匀。
\end{enumerate}

%% answer:
\jdanswer{
\begin{enumerate}
	\item
	BC
	\item
	（Ⅰ）$OA>0.20\Um$；（Ⅱ）$OA>0.25\Um$。
\end{enumerate}
}

%% memo:
\memoanswer{
（1）由皮带传动，从动轮与主动轮的转动方向相反，故从动轮做逆时针转动，B 正确；轮缘处线速度的大小相等，有 $n_{2}r_{2}=nr_{1}$，故从动轮的转速 $n_{2}=\frac{r_{1}}{r_{2}}n$，C 正确。

（2）（Ⅰ）当木条 A 端刚刚离开斜面时，受力情况如图。

\onepicture[w=5.0cm]{figs/30_an.png}

设斜面倾角为 $\theta$，根据力矩平衡条件，若满足 $Mg\cdot OA\cos\theta>mg\cdot OB\cos\theta$（①），木条就不会脱离斜面。根据题意有 $OA+OB=l$（②）。联立①②并代入已知条件得 $OA>0.20\Um$（③）。

（Ⅱ）设 $G$ 为木条重心，由题意可知 $AG=\frac{1}{2}l$（④）。当木条 A 端刚刚离开斜面时，若满足 $Mg\cdot OA\cos\theta>mg\cdot OB\cos\theta+mg\cdot OG\cos\theta$（⑤），木条就不会脱离斜面。联立②④⑤并代入已知条件得 $OA>0.25\Um$（⑥）。
}

\item
%% number: 31
%% paperName: 2008年宁夏理综第 31 题 9 分
%% typeId: 6
%% type: 计算题
%% chapter: 气体性质
%% point: 玻意耳定律
%% method: 无
%% score: 9
%% degree: 650
%% duplicateId: 0
%% body:
\begin{enumerate}
	\item
	（6 分）如图所示，由导热材料制成的气缸和活塞将一定质量的理想气体封闭在气缸内，活塞与气缸壁之间无摩擦，活塞上方存有少量液体。将一细管插入液体，由于虹吸现象，活塞上方液体逐渐流出。在此过程中，大气压强与外界的温度保持不变。关于这一过程，下列说法正确的是 \tkanswer[1.0cm]{}。（填入选项前的字母，有填错的不得分）
	
	\onepicture[w=4.0cm]{figs/31a.png}
	
	（A）气体分子的平均动能逐渐增大
	
	（B）单位时间气体分子对活塞撞击的次数增多
	
	（C）单位时间气体分子对活塞的冲量保持不变
	
	（D）气体对外界做功等于气体从外界吸收的热量
	
	\item
	（9 分）一定质量的理想气体被活塞封闭在可导热的气缸内，活塞相对于底部的高度为 $h$，可沿气缸无摩擦地滑动。取一小盒沙子缓慢地倒在活塞的上表面上。沙子倒完时，活塞下降了 $\frac{1}{4}h$。再取相同质量的一小盒沙子缓慢地倒在活塞的上表面上。外界大气的压强和温度始终保持不变，求此次沙子倒完时活塞距气缸底部的高度。
	
	\onepicture[w=3.0cm, align=right]{figs/31b.png}
\end{enumerate}

%% answer:
\jdanswer{
\begin{enumerate}
	\item
	D
	\item
	$h'=\frac{3}{5}h$
\end{enumerate}
}

%% memo:
\memoanswer{
（1）气体压强等于大气压与活塞、液体重力产生的压强之和；液体逐渐流出，气体压强逐渐减小，气体体积增大，气体对外界做功；又因外界温度不变，气体内能不变，由热力学第一定律可知气体对外界做的功等于从外界吸收的热量，D 正确。

（2）设大气和活塞对气体的总压强为 $p_{0}$，加一小盒沙子对气体产生的压强为 $p$，由玻意耳定律得 $p_{0}h=(p_{0}+p)\left(h-\frac{1}{4}h\right)$（①），解得 $p=\frac{1}{3}p_{0}$（②）。再加一小盒沙子后，气体的压强变为 $p_{0}+2p$，设第二次加沙子后活塞的高度为 $h'$，则 $p_{0}h=(p_{0}+2p)h'$（③）。联立②③式解得 $h'=\frac{3}{5}h$。
}

\item
%% number: 32
%% paperName: 2008年宁夏理综第 32 题 9 分
%% typeId: 6
%% type: 计算题
%% chapter: 光学
%% point: 光的折射
%% method: 无
%% score: 9
%% degree: 650
%% duplicateId: 0
%% body:
\begin{enumerate}
	\item
	（6 分）下列关于简谐振动和简谐机械波的说法正确的是 \tkanswer[1.0cm]{}。（填入选项前的字母，有填错的不得分）
	
	（A）弹簧振子的周期与振幅有关
	
	（B）横波在介质中的传播速度由介质本身的性质决定
	
	（C）在波传播方向上的某个质点的振动速度就是波的传播速度
	
	（D）单位时间内经过媒质中一点的完全波的个数就是这列简谐波的频率
	
	\item
	（9 分）一半径为 $R$ 的 $\frac{1}{4}$ 球体放置在水平面上，球体由折射率为 $\sqrt{3}$ 的透明材料制成。现有一束位于过球心 O 的竖直平面内的光线，平行于桌面射到球体表面上，折射入球体后再从竖直表面射出，如图所示。已知入射光线与桌面的距离为 $\frac{\sqrt{3}}{2}R$。求出射角 $\theta$。
	
	\onepicture[w=5.0cm, align=right]{figs/32.png}
\end{enumerate}

%% answer:
\jdanswer{
\begin{enumerate}
	\item
	BD
	\item
	$\theta=60^{\circ}$
\end{enumerate}
}

%% memo:
\memoanswer{
（1）弹簧振子的周期由系统本身决定，与振幅无关，A 错；横波在介质中的传播速度由介质本身的性质决定，B 对；波传播方向上某个质点的振动速度与波的传播速度是两回事，C 错；单位时间内经过媒质中一点的完全波的个数就是这列简谐波的频率，D 对。故选 BD。

（2）设入射光线与 $\frac{1}{4}$ 球体的交点为 C，连接 OC，OC 即为入射点的法线，图中的角 $\alpha$ 为入射角。

\onepicture[w=5.5cm]{figs/32_an.png}

过 C 点作球体水平表面的垂线，垂足为 B，则 $\angle COB=\alpha$。由 $\triangle OBC$ 知 $\sin\alpha=\frac{\sqrt{3}}{2}$（①）。设光线在 C 点的折射角为 $\beta$，由折射定律得 $\frac{\sin\alpha}{\sin\beta}=\sqrt{3}$（②）。由①②式得 $\beta=30^{\circ}$（③）。由几何关系知，光线在球体的竖直表面上的入射角 $\gamma=30^{\circ}$（④）。由折射定律得 $\frac{\sin\gamma}{\sin\theta}=\frac{1}{\sqrt{3}}$（⑤）。因此 $\sin\theta=\frac{\sqrt{3}}{2}$，解得 $\theta=60^{\circ}$。
}

\item
%% number: 33
%% paperName: 2008年宁夏理综第 33 题 9 分
%% typeId: 6
%% type: 计算题
%% chapter: 运动和力
%% point: 动量守恒定律
%% method: 无
%% score: 9
%% degree: 650
%% duplicateId: 0
%% body:
\begin{enumerate}
	\item
	（6 分）天然放射性元素 $\ce{^{239}_{94}Pu}$ 经过 \tkanswer[1.0cm]{} 次 $\alpha$ 衰变和 \tkanswer[1.0cm]{} 次 $\beta$ 衰变，最后变成铅的同位素 \tkanswer[2.0cm]{}。（填入铅的三种同位素 $\ce{^{206}_{82}Pb}$、$\ce{^{207}_{82}Pb}$、$\ce{^{208}_{82}Pb}$ 中的一种）
	
	\item
	（9 分）某同学利用如图所示的装置验证动量守恒定律。图中两摆摆长相同，悬挂于同一高度，A、B 两摆球均很小，质量之比为 $1:2$。当两摆均处于自由静止状态时，其侧面刚好接触。向右上方拉动 B 球使其摆线伸直并与竖直方向成 $45^{\circ}$ 角，然后将其由静止释放。结果观察到两摆球粘在一起摆动，且最大摆角成 $30^{\circ}$。若本实验允许的最大误差为 $\pm4\%$，此实验是否成功地验证了动量守恒定律？
	
	\onepicture[w=2.5cm, align=right]{figs/33.png}
\end{enumerate}

%% answer:
\jdanswer{
\begin{enumerate}
	\item
	8，4，$\ce{^{207}_{82}Pb}$
	\item
	是，此实验在规定的范围内验证了动量守恒定律。
\end{enumerate}
}

%% memo:
\memoanswer{
（1）每发生一次 $\alpha$ 衰变，质量数减少 4、电荷数减少 2；每发生一次 $\beta$ 衰变，质量数不变、电荷数增加 1。由 $\ce{^{239}_{94}Pu}$ 变为 $\ce{^{207}_{82}Pb}$，质量数减少 32，故 $\alpha$ 衰变次数为 $\frac{239-207}{4}=8$；电荷数应减少 $2\times8=16$，实际减少 $94-82=12$，故 $\beta$ 衰变次数为 $16-12=4$。

（2）设摆球 A、B 的质量分别为 $m_{A}$、$m_{B}$，摆长为 $l$，B 球的初始高度为 $h_{1}$，碰撞前 B 球的速度为 $v_{B}$。在不考虑摆线质量的情况下，根据题意及机械能守恒定律得 $h_{1}=l(1-\cos45^{\circ})$（①），$\frac{1}{2}m_{B}v_{B}^{2}=m_{B}gh_{1}$（②）。

设碰撞前、后两摆球的总动量的大小分别为 $p_{1}$、$p_{2}$，有 $p_{1}=m_{B}v_{B}$（③）。联立①②③式得 $p_{1}=m_{B}\sqrt{2gl(1-\cos45^{\circ})}$（④）。同理可得 $p_{2}=(m_{A}+m_{B})\sqrt{2gl(1-\cos30^{\circ})}$（⑤）。联立④⑤式得 $\frac{p_{2}}{p_{1}}=\frac{m_{A}+m_{B}}{m_{B}}\sqrt{\frac{1-\cos30^{\circ}}{1-\cos45^{\circ}}}$（⑥）。代入已知条件得 $\left(\frac{p_{2}}{p_{1}}\right)^{2}=1.03$（⑦）。由此可以推出 $\frac{p_{2}-p_{1}}{p_{1}}\le4\%$（⑧）。所以，此实验在规定的范围内验证了动量守恒定律。
}

\end{enumerate}

\end{document}
